\section{Step 2 -- Signature axes via maximization principle}

% ---------------------------------------------------------------
\subsection*{Ziel und Grundidee}

Step~2 macht die in Step~1 gefundenen initialen Achsen unabh\"{a}ngig
vom zuf\"{a}lligen Startvertreter, indem sukzessive weitere
repr\"{a}sentative Gene und Samples gesammelt und zu akkumulierten
Achsen kombiniert werden, so lange bis diese gegen ein lokales
Maximum des Signaturfunktionals $\xi$ konvergieren.

Die akkumulierten Achsen nach $l$ Repr\"{a}sentanten werden mit
$(|\mathbf{b}^g_l\rangle, |\mathbf{b}^s_l\rangle)$ bezeichnet.
F\"{u}r $l = 1$ gilt per Definition:
\begin{equation}
    |\mathbf{b}^g_1\rangle \;\coloneqq\; |a^g\rangle,
    \qquad
    |\mathbf{b}^s_1\rangle \;\coloneqq\; |a^s\rangle,
\end{equation}
mit den Korrelationen $|r^g_1\rangle,\,|r^s_1\rangle$,
Gewichten $|w^g_1\rangle,\,|w^s_1\rangle$,
$p$-Werten $|p^g_1\rangle,\,|p^s_1\rangle$
und Signaturfunktional $\mathcal{E}_1$ direkt aus Step~1.

% ---------------------------------------------------------------
\subsection*{Iterationsschleife \"{u}ber $l = 1, 2, 3, \ldots$}

% ---------------------------------------------------------------
\paragraph{1. Sortierung der Kandidaten.}

Alle Gene und Samples werden nach ihren absoluten Korrelationen
zur aktuellen akkumulierten Achse $|\mathbf{b}^g_l\rangle$
bzw.\ $|\mathbf{b}^s_l\rangle$ absteigend sortiert:
\begin{align}
    \bigl|\langle e^g_{i_1} \mid r^g_l \rangle\bigr|
    &\;\geq\;
    \bigl|\langle e^g_{i_2} \mid r^g_l \rangle\bigr|
    \;\geq\; \cdots,
    \\[4pt]
    \bigl|\langle e^s_{j_1} \mid r^s_l \rangle\bigr|
    &\;\geq\;
    \bigl|\langle e^s_{j_2} \mid r^s_l \rangle\bigr|
    \;\geq\; \cdots
\end{align}
Nur die Top-Kandidaten dieser sortierten Liste werden als
m\"{o}gliche n\"{a}chste Repr\"{a}sentanten getestet.

% ---------------------------------------------------------------
\paragraph{2. Auswahl des n\"{a}chsten Repr\"{a}sentanten.}

F\"{u}r jeden Top-Kandidaten werden hypothetisch die akkumulierten
Achsen $|\mathbf{b}^g_{l+1}\rangle$ und $|\mathbf{b}^s_{l+1}\rangle$
nach Gl.~(9) und Gl.~(10) berechnet sowie die zugeh\"{o}rigen
Korrelationen $|r^g_{l+1}\rangle$ und $|r^s_{l+1}\rangle$,
$p$-Werte $|p^g_{l+1}\rangle$ und $|p^s_{l+1}\rangle$
und das Signaturfunktional $\xi_{l+1}$.
Ausgew\"{a}hlt wird derjenige Kandidat, der $\mathcal{E}_{l+1}$
maximiert, und dieser wird zum Repr\"{a}sentanten $l+1$.

% ---------------------------------------------------------------
\paragraph{3. Akkumulation der Achsen (Gl.~9 und Gl.~10).}

Die neue akkumulierte Sample-Achse ist das
korrelationsgewichtete arithmetische Mittel \"{u}ber alle bisherigen
Einzel-Achsen $|a^s_{l'}\rangle$, $l' = 1, \ldots, l+1$:
\begin{equation}
    |\mathbf{b}^s_{l+1}\rangle
    \;\equiv\;
    \frac{\displaystyle
          \sum_{l'=1}^{l+1}
          \bigl[\langle e^s_{l'} \mid \mathbf{b}^s_{l'} \rangle\bigr]_{|w^s_{l'}\rangle}
          \cdot |a^s_{l'}\rangle}
         {\displaystyle
          \sum_{l'=1}^{l+1}
          \Bigl|\bigl[\langle e^s_{l'} \mid \mathbf{b}^s_{l'} \rangle\bigr]_{|w^s_{l'}\rangle}\Bigr|}.
\end{equation}
Analog gilt f\"{u}r die Gen-Achse:
\begin{equation}
    |\mathbf{b}^g_{l+1}\rangle
    \;\equiv\;
    \frac{\displaystyle
          \sum_{l'=1}^{l+1}
          \bigl[\langle e^g_{l'} \mid \mathbf{b}^g_{l'} \rangle\bigr]_{|w^g_{l'}\rangle}
          \cdot |a^g_{l'}\rangle}
         {\displaystyle
          \sum_{l'=1}^{l+1}
          \Bigl|\bigl[\langle e^g_{l'} \mid \mathbf{b}^g_{l'} \rangle\bigr]_{|w^g_{l'}\rangle}\Bigr|}.
\end{equation}
Das Gewicht jedes Repr\"{a}sentanten $l'$ ist seine Korrelation
zur bisherigen akkumulierten Achse.
Repr\"{a}sentanten, die schlecht zur akkumulierten Richtung passen,
tragen kaum bei, w\"{a}hrend gut passende Repr\"{a}sentanten vollen
Einfluss behalten.

% ---------------------------------------------------------------
\paragraph{4. Konvergenzpr\"{u}fung (Supp.~Note~6).}

Nach jeder Aktualisierung wird gepr\"{u}ft, ob beide folgenden
Kriterien gleichzeitig erf\"{u}llt sind.

Das erste Kriterium fordert geometrische Achsenkonvergenz.
Es wird gemessen, wie stark sich die akkumulierten Achsen noch
von Iteration zu Iteration ver\"{a}ndern, als eins minus der
Korrelation zur vorherigen Achse:
\begin{equation}
    \delta^g \;\equiv\;
    1 - \bigl[\langle \mathbf{b}^g_l \mid \mathbf{b}^g_{l+1}
               \rangle\bigr]_{|w^g_{l+1}\rangle},
    \qquad
    \delta^s \;\equiv\;
    1 - \bigl[\langle \mathbf{b}^s_l \mid \mathbf{b}^s_{l+1}
               \rangle\bigr]_{|w^s_{l+1}\rangle}.
\end{equation}
Konvergenz liegt vor, wenn gilt:
\begin{equation}
    \frac{\delta^g + \delta^s}{2} \;<\; 10^{-3}.
\end{equation}
Dies entspricht einer Achsen\"{a}hnlichkeit von mehr als $0{,}999$.
Die Konvergenz ist garantiert, da jeder neue Repr\"{a}sentant die
akkumulierte Achse nur noch mit einem maximalen Gewicht von $1/l$
ver\"{a}ndern kann.

Das zweite Kriterium fordert eine Mindestanzahl an akkumulierten
Repr\"{a}sentanten:
\begin{equation}
    l \;\geq\; \min\!\bigl(15,\; 0{,}2 \cdot m_k\bigr)
    \qquad\text{und}\qquad
    l \;\geq\; \min\!\bigl(15,\; 0{,}2 \cdot n_k\bigr).
\end{equation}
Dies sch\"{u}tzt vor zuf\"{a}lliger Fr\"{u}hkonvergenz, etwa wenn
zwei nahezu identische Startgene vorliegen, sodass die Mindestanzahl
sicherstellt, dass die Achse auf einer ausreichend breiten Basis
von Repr\"{a}sentanten beruht.

Sind beide Kriterien erf\"{u}llt, wird die Iteration abgebrochen
und $\hat{l}$ ist erreicht.
Andernfalls wird mit der Suche nach dem n\"{a}chsten Repr\"{a}sentanten
fortgefahren, beginnend wieder bei Schritt~1.

% ---------------------------------------------------------------
\subsection*{Ergebnis von Step~2}

Die finalen akkumulierten Achsen bei $\hat{l}$,
\begin{equation}
    \bigl(|\mathbf{b}^g_{\hat{l}}\rangle,\;
           |\mathbf{b}^s_{\hat{l}}\rangle\bigr),
\end{equation}
stellen die lineare Approximation der zugrundeliegenden
Signaturinteraktion dar und werden an Step~3 \"{u}bergeben.